\chapter{Integral de línea de un campo vectorial y teorema de Green}

La integral de línea de un campo escalar, estudiada en la lección anterior,
permitía calcular la masa de un alambre o el área de una superficie curva,
acumulando una función a lo largo de una trayectoria. Ahora se aborda un
problema más sutil: la integral de línea de un \emph{campo vectorial}, que
mide la tendencia del campo a empujar o girar a lo largo de la curva. En
contextos físicos, esta integral se interpreta como el \emph{trabajo}
realizado por una fuerza al desplazar una partícula a lo largo de una
trayectoria, o como la \emph{circulación} de un fluido alrededor de un
obstáculo. A diferencia del caso escalar, el resultado de esta integral
depende críticamente de la orientación de la curva y, en general, del
camino seguido.

El objetivo de esta lección es doble. Por un lado, se estudian las
condiciones bajo las cuales la integral de línea de un campo vectorial
\emph{no depende} del camino, sino únicamente de los puntos inicial y
final, lo que conduce al \textbf{teorema fundamental de las integrales
de línea} y a la noción de \textbf{campo conservativo}. Por otro lado,
cuando la curva es cerrada, se establece una conexión profunda entre la
circulación a lo largo de la frontera y el comportamiento interno del
campo en la región encerrada: esta conexión es el \textbf{teorema de
Green}, que relaciona una integral de línea con una integral doble del
rotacional. Este teorema es un caso particular en el plano de resultados
más generales (Stokes, divergencia) que se estudiarán en la lección
siguiente, y constituye una de las piedras angulares del cálculo
vectorial.

%%==========================================================
\section{Integral de línea de un campo vectorial}
%%==========================================================

La definición de integral de línea de un campo vectorial se construye
siguiendo el mismo principio de Riemann que ya se ha usado para
integrales de funciones escalares: se divide la curva en segmentos
pequeños, se multiplica la componente del campo en la dirección del
desplazamiento por la longitud de cada segmento, y se suman todas las
contribuciones. Al tomar el límite cuando los segmentos se hacen
infinitamente pequeños, se obtiene la integral.

\begin{definition}[Integral de línea de un campo vectorial]
Sea $\mathbf{F}$ un campo vectorial continuo definido en una región que
contiene una curva suave por partes $C$, parametrizada por
\[
\mathbf{r}(t) = \langle x(t), y(t), z(t) \rangle, \qquad a \le t \le b,
\]
con $\mathbf{r}'$ continua y no nula en cada subintervalo. La
\textbf{integral de línea de $\mathbf{F}$ a lo largo de $C$} se define
como
\[
\int_C \mathbf{F} \cdot d\mathbf{r}
= \int_a^b \mathbf{F}(\mathbf{r}(t)) \cdot \mathbf{r}'(t) \, dt.
\]
Si $\mathbf{F} = \langle M, N, P \rangle$ y $d\mathbf{r} = \langle dx,
dy, dz \rangle$, también se escribe
\[
\int_C M\,dx + N\,dy + P\,dz.
\]
La integral es independiente de la parametrización elegida, siempre que
se mantenga la misma orientación de $C$.
\end{definition}

La integral puede interpretarse como la suma de las proyecciones del
campo sobre la dirección tangente a la curva: si $\mathbf{T}(t)$ es el
vector tangente unitario y $ds$ el elemento de longitud de arco,
entonces
\[
\int_C \mathbf{F} \cdot d\mathbf{r}
= \int_C \mathbf{F} \cdot \mathbf{T} \, ds.
\]

\begin{rem}
Dos diferencias fundamentales con la integral de línea de un campo
escalar:
\begin{itemize}
    \item La integral de línea de un campo vectorial \emph{cambia de
    signo} si se invierte la orientación de la curva: $\int_{-C}
    \mathbf{F}\cdot d\mathbf{r} = -\int_C \mathbf{F}\cdot d\mathbf{r}$.
    La integral de un campo escalar no depende de la orientación.
    \item La integral de línea de un campo vectorial mide una
    \emph{circulación} o \emph{trabajo}; es sensible a la dirección en
    que se recorre la curva.
\end{itemize}
\end{rem}

\begin{pasos}[Cálculo de una integral de línea de un campo vectorial]
Para evaluar $\displaystyle \int_C \mathbf{F} \cdot d\mathbf{r}$:
\begin{enumerate}
    \item \textbf{Parametrizar la curva.} Hallar una parametrización
    suave por partes $\mathbf{r}(t)$, $a \le t \le b$, que recorra $C$
    en la orientación dada.
    \item \textbf{Calcular el vector tangente.} Derivar $\mathbf{r}' (t)$.
    \item \textbf{Evaluar el campo sobre la curva.} Formar
    $\mathbf{F}(\mathbf{r}(t))$.
    \item \textbf{Calcular el producto punto.} Obtener
    $\mathbf{F}(\mathbf{r}(t)) \cdot \mathbf{r}'(t)$, que es una
    función escalar de $t$.
    \item \textbf{Integrar.} Calcular $\displaystyle \int_a^b
    \mathbf{F}(\mathbf{r}(t)) \cdot \mathbf{r}'(t) \, dt$.
\end{enumerate}
\end{pasos}

%%==========================================================
\begin{example}[Trabajo de una fuerza variable]
Calcular el trabajo realizado por el campo de fuerzas
$\mathbf{F}(x,y) = \langle y, x^2 \rangle$ al mover una partícula desde
$(0,0)$ hasta $(1,1)$ a lo largo de la parábola $y=x^2$.
\end{example}

\begin{myproof}
\textbf{Paso 1. Parametrizar la curva.}
La parábola $y=x^2$ se parametriza como
\[
\mathbf{r}(t) = \langle t, t^2 \rangle, \qquad 0 \le t \le 1,
\]
que recorre la curva desde $(0,0)$ hasta $(1,1)$ en la orientación
correcta.

\textbf{Paso 2. Calcular el vector tangente.}
Derivando:
\[
\mathbf{r}'(t) = \langle 1, 2t \rangle.
\]

\textbf{Paso 3. Evaluar el campo sobre la curva.}
Sustituyendo $x=t$, $y=t^2$ en $\mathbf{F}$:
\[
\mathbf{F}(\mathbf{r}(t)) = \langle t^2, t^2 \rangle.
\]

\textbf{Paso 4. Calcular el producto punto.}
\[
\mathbf{F}(\mathbf{r}(t)) \cdot \mathbf{r}'(t)
= \langle t^2, t^2 \rangle \cdot \langle 1, 2t \rangle
= t^2 + 2t^3.
\]

\textbf{Paso 5. Integrar.}
\[
\int_C \mathbf{F} \cdot d\mathbf{r}
= \int_0^1 (t^2 + 2t^3) \, dt
= \left[ \frac{t^3}{3} + \frac{t^4}{2} \right]_0^1
= \frac{1}{3} + \frac{1}{2} = \frac{5}{6}.
\]
Por tanto, el trabajo realizado es $\boxed{\dfrac{5}{6}}$ unidades de
trabajo (en las unidades correspondientes a $\mathbf{F}$ y la
parametrización). $\square$
\end{myproof}

\begin{rem}
Si se hubiera recorrido la misma curva pero en sentido opuesto, la
parametrización $\mathbf{r}(t)=\langle 1-t, (1-t)^2 \rangle$,
$0\le t\le 1$, habría producido el valor $-\frac{5}{6}$. Esta
dependencia de la orientación es característica de la integral de línea
de campos vectoriales y refleja la diferencia entre el trabajo realizado
por una fuerza y el trabajo realizado contra ella.
\end{rem}

%%==========================================================
\section{Independencia de la trayectoria y campos conservativos}
%%==========================================================

En el ejemplo anterior, si se hubiera elegido otro camino entre los
mismos puntos, por ejemplo la línea recta $\mathbf{r}(t)=\langle t,t
\rangle$, el resultado habría sido
\[
\int_C \mathbf{F}\cdot d\mathbf{r}
= \int_0^1 \langle t, t^2 \rangle \cdot \langle 1,1 \rangle \, dt
= \int_0^1 (t+t^2) dt = \frac{5}{6}.
\]
En este caso particular, el resultado coincidió con el de la parábola.
Sin embargo, \emph{no siempre ocurre así}. La pregunta fundamental es:
¿bajo qué condiciones la integral de línea de un campo vectorial
\emph{no depende} del camino elegido, sino únicamente de los puntos
extremos?

La respuesta conduce a la noción de \textbf{campo conservativo}, que
es el análogo en varias variables de la antiderivada en una variable.

\begin{definition}[Campo conservativo]
Un campo vectorial $\mathbf{F}$ definido en una región $D$ se dice
\textbf{conservativo} si existe una función escalar $f$ (llamada
\textbf{función potencial}) tal que
\[
\mathbf{F} = \nabla f
\]
en todo punto de $D$. Es decir, $\mathbf{F}$ es el gradiente de $f$.
\end{definition}

\begin{theorem}[Independencia de la trayectoria]
Sea $\mathbf{F}$ un campo vectorial continuo en una región abierta y
conexa $D$. La integral de línea
\[
\int_C \mathbf{F} \cdot d\mathbf{r}
\]
es \textbf{independiente de la trayectoria} en $D$ (es decir, su valor
solo depende de los puntos inicial y final, no de la curva que los une)
si y solo si $\mathbf{F}$ es conservativo en $D$.
\end{theorem}

\begin{proof}
(Esquema de la demostración.) Si $\mathbf{F}=\nabla f$, entonces,
usando la regla de la cadena,
\[
\mathbf{F}(\mathbf{r}(t)) \cdot \mathbf{r}'(t)
= \nabla f(\mathbf{r}(t)) \cdot \mathbf{r}'(t)
= \frac{d}{dt} f(\mathbf{r}(t)).
\]
Por el teorema fundamental del cálculo,
\[
\int_C \mathbf{F} \cdot d\mathbf{r}
= \int_a^b \frac{d}{dt} f(\mathbf{r}(t)) \, dt
= f(\mathbf{r}(b)) - f(\mathbf{r}(a)),
\]
que solo depende de los extremos.

Recíprocamente, si la integral es independiente de la trayectoria, se
puede definir $f(P)$ como la integral de $\mathbf{F}$ desde un punto
fijo $P_0$ hasta $P$, y se demuestra que $\nabla f = \mathbf{F}$.
La demostración formal requiere verificar que las derivadas parciales
de $f$ coinciden con las componentes de $\mathbf{F}$.
\end{proof}

\begin{theorem}[Teorema fundamental de las integrales de línea]
Sea $C$ una curva suave por partes que une los puntos $A$ y $B$,
contenida en una región $D$. Si $\mathbf{F} = \nabla f$ es un campo
gradiente continuo en $D$, entonces
\[
\int_C \mathbf{F} \cdot d\mathbf{r} = f(B) - f(A).
\]
En particular, si $C$ es cerrada ($A=B$), entonces
\[
\oint_C \mathbf{F} \cdot d\mathbf{r} = 0.
\]
\end{theorem}

\begin{rem}
El teorema fundamental generaliza la regla de Barrow para integrales
definidas: la integral del gradiente de $f$ a lo largo de una curva
es simplemente la diferencia de $f$ en los extremos. Esto simplifica
radicalmente el cálculo de integrales de línea: en lugar de parametrizar
la curva, basta con encontrar la función potencial y evaluarla en los
extremos. La dificultad se traslada entonces a determinar si el campo
es conservativo y, en caso afirmativo, hallar $f$.
\end{rem}

\begin{pasos}[Determinación de un campo conservativo y su potencial]
Para verificar si un campo $\mathbf{F}=\langle M,N,P \rangle$ es
conservativo y encontrar su función potencial:
\begin{enumerate}
    \item \textbf{Verificar el dominio.} La región debe ser abierta y
    conexa. En $\mathbb{R}^2$, si el dominio es simplemente conexo
    (sin agujeros), basta comprobar la condición de simetría de las
    derivadas cruzadas:
    \[
    \frac{\partial M}{\partial y} = \frac{\partial N}{\partial x}.
    \]
    En $\mathbb{R}^3$, se requiere $\nabla \times \mathbf{F} = \mathbf{0}$
    (el rotacional es cero).
    \item \textbf{Integrar una componente.} Si $\mathbf{F} = \nabla f$,
    entonces $\partial f/\partial x = M$. Integrar $M$ respecto a $x$:
    \[
    f(x,y,z) = \int M\,dx + g(y,z),
    \]
    donde $g$ es una función que depende solo de las variables
    restantes.
    \item \textbf{Derivar y ajustar.} Derivar la expresión anterior
    respecto a $y$ e igualar a $N$ para determinar $\partial g/\partial y$.
    Repetir para $P$ si estamos en $\mathbb{R}^3$.
    \item \textbf{Construir $f$.} Una vez determinadas todas las
    constantes de integración, escribir $f$ explícitamente.
\end{enumerate}
\end{pasos}

\begin{rem}
La condición de simetría de las derivadas cruzadas ($\partial M/\partial
y = \partial N/\partial x$ en $\mathbb{R}^2$) es necesaria para que el
campo sea conservativo, pero solo es suficiente si el dominio es
\emph{simplemente conexo}. En un dominio con agujeros, puede fallar:
por ejemplo, el campo $\mathbf{F}(x,y) = \langle -y/(x^2+y^2),
x/(x^2+y^2) \rangle$ satisface la condición de simetría en su dominio
$\mathbb{R}^2\setminus\{(0,0)\}$, pero no es conservativo, ya que su
integral alrededor de la circunferencia unidad es $2\pi \neq 0$.
\end{rem}

%%==========================================================
\begin{example}[Aplicación del teorema fundamental]
Verificar que el campo $\mathbf{F}(x,y) = \langle 2xy, x^2 + 2y \rangle$
es conservativo, encontrar su función potencial, y calcular
$\int_C \mathbf{F}\cdot d\mathbf{r}$ a lo largo de cualquier curva que
una $(0,1)$ con $(1,2)$.
\end{example}

\begin{myproof}
\textbf{Paso 1. Verificar la condición de simetría.}
Aquí $M=2xy$, $N=x^2+2y$. Se calcula:
\[
\frac{\partial M}{\partial y} = 2x, \qquad
\frac{\partial N}{\partial x} = 2x.
\]
Como son iguales en todo $\mathbb{R}^2$ (que es simplemente conexo),
el campo es conservativo.

\textbf{Paso 2. Integrar una componente.}
Se integra $M$ respecto a $x$:
\[
f(x,y) = \int 2xy \, dx = x^2 y + g(y),
\]
donde $g$ es una función de $y$ (la constante de integración en $x$).

\textbf{Paso 3. Derivar y ajustar.}
Derivando $f$ respecto a $y$ e igualando a $N$:
\[
\frac{\partial f}{\partial y} = x^2 + g'(y) = x^2 + 2y.
\]
Por tanto, $g'(y) = 2y$, de donde $g(y) = y^2 + C$. Tomando $C=0$,
\[
f(x,y) = x^2 y + y^2.
\]

\textbf{Paso 4. Construir $f$.}
La función potencial es $f(x,y)=x^2 y + y^2$.

\textbf{Paso 5. Aplicar el teorema fundamental.}
Para cualquier curva que una $A=(0,1)$ y $B=(1,2)$:
\[
\int_C \mathbf{F}\cdot d\mathbf{r}
= f(1,2) - f(0,1)
= (1^2\cdot 2 + 2^2) - (0^2\cdot 1 + 1^2)
= (2+4) - 1 = 5.
\]
La integral vale $\boxed{5}$, independientemente de la trayectoria.
$\square$
\end{myproof}

\begin{rem}
Nótese que la función potencial no es única: si se añade una constante
$C$, la diferencia $f(B)-f(A)$ no se altera. Por tanto, la constante
de integración es irrelevante para el cálculo de la integral de línea.
\end{rem}

%%==========================================================
\section{Teorema de Green}
%%==========================================================

El teorema de Green establece una conexión entre una integral de línea
sobre una curva cerrada simple y una integral doble sobre la región
encerrada por esa curva. Es el primer resultado que relaciona el
comportamiento de un campo \emph{en el interior} de una región con su
comportamiento \emph{en la frontera}, y constituye el germen de los
teoremas de Stokes y de la divergencia que se estudiarán en la lección
siguiente.

\begin{definition}[Curva cerrada simple y orientación positiva]
Una curva $C$ en el plano se dice \textbf{cerrada} si su punto inicial
coincide con su punto final, y \textbf{simple} si no se intersecta a sí
misma (excepto en los extremos). Su \textbf{orientación positiva} es
aquella en la que la región encerrada queda siempre a la izquierda al
recorrer la curva (sentido antihorario para una curva que encierra una
región convexa).
\end{definition}

\begin{theorem}[Teorema de Green]
Sea $C$ una curva suave por partes, cerrada simple, orientada
positivamente en el plano, y sea $R$ la región encerrada por $C$.
Si $M(x,y)$ y $N(x,y)$ son funciones con derivadas parciales continuas
en una región abierta que contiene a $R$, entonces
\[
\oint_C M\,dx + N\,dy
= \iint_R \left( \frac{\partial N}{\partial x} - \frac{\partial M}{\partial y} \right) dA.
\]
Equivalentemente, si $\mathbf{F} = \langle M, N \rangle$,
\[
\oint_{\partial R} \mathbf{F} \cdot d\mathbf{r}
= \iint_R (\nabla \times \mathbf{F}) \cdot \mathbf{k} \, dA,
\]
donde $\nabla \times \mathbf{F} = \left( \frac{\partial N}{\partial x}
- \frac{\partial M}{\partial y} \right) \mathbf{k}$ es la componente
$z$ del rotacional.
\end{theorem}

\begin{proof}[Demostración pedagógica para regiones simples]
Basta demostrar el teorema para una región que se pueda describir
simultáneamente como una región de tipo I y de tipo II. Supóngase que
$R = \{(x,y) \mid a \le x \le b,\ g_1(x) \le y \le g_2(x)\}$. La
frontera se compone de cuatro curvas. Para la integral de $M\,dx$,
se usa el teorema fundamental del cálculo:
\[
\oint_C M\,dx
= \int_{C_1} M\,dx + \int_{C_2} M\,dx
= \int_a^b M(x,g_1(x))\,dx - \int_a^b M(x,g_2(x))\,dx,
\]
donde $C_1$ es la frontera inferior y $C_2$ la superior (orientadas
correctamente). Por el teorema fundamental,
\[
\int_{g_1(x)}^{g_2(x)} \frac{\partial M}{\partial y}\,dy
= M(x,g_2(x)) - M(x,g_1(x)).
\]
Por tanto,
\[
\oint_C M\,dx = - \int_a^b \int_{g_1(x)}^{g_2(x)}
\frac{\partial M}{\partial y}\,dy\,dx
= -\iint_R \frac{\partial M}{\partial y}\,dA.
\]
De manera análoga, expresando $R$ como región de tipo II
($c \le y \le d,\ h_1(y) \le x \le h_2(y)$), se obtiene
\[
\oint_C N\,dy = \iint_R \frac{\partial N}{\partial x}\,dA.
\]
Sumando ambas igualdades se obtiene la fórmula de Green.
\end{proof}

\begin{rem}
La demostración anterior supone que la región es simple (se puede
describir como tipo I y tipo II). Para regiones más generales (con
agujeros o fronteras más complejas), se descomponen en subregiones
simples y se aplica Green en cada una, sumando los resultados. Las
integrales de línea sobre los cortes internos se cancelan, dejando
únicamente la integral sobre la frontera exterior.
\end{rem}

\begin{pasos}[Aplicación del teorema de Green]
Para usar el teorema de Green:
\begin{enumerate}
    \item \textbf{Identificar $M$ y $N$.} Escribir el campo vectorial
    como $\mathbf{F}=\langle M,N \rangle$ y la integral de línea como
    $\oint_C M\,dx + N\,dy$.
    \item \textbf{Calcular el integrando de la integral doble.}
    Hallar $\displaystyle \frac{\partial N}{\partial x} -
    \frac{\partial M}{\partial y}$.
    \item \textbf{Describir la región $R$.} Identificar la región
    encerrada por $C$ y determinar sus límites de integración.
    \item \textbf{Evaluar la integral doble.} Calcular
    $\displaystyle \iint_R \left( \frac{\partial N}{\partial x} -
    \frac{\partial M}{\partial y} \right) dA$.
\end{enumerate}
\end{pasos}

%%==========================================================
\begin{example}[Verificación de Green]
Verificar el teorema de Green para el campo $\mathbf{F}(x,y) =
\langle -y, x \rangle$ y la región $R$ encerrada por el círculo
$x^2+y^2=4$.
\end{example}

\begin{myproof}
\textbf{Paso 1. Identificar $M$ y $N$.}
Aquí $M=-y$, $N=x$.

\textbf{Paso 2. Calcular el integrando de la integral doble.}
\[
\frac{\partial N}{\partial x} - \frac{\partial M}{\partial y}
= \frac{\partial}{\partial x}(x) - \frac{\partial}{\partial y}(-y)
= 1 - (-1) = 2.
\]

\textbf{Paso 3. Describir la región $R$.}
$R$ es el disco de radio $2$ centrado en el origen:
$R = \{(x,y) \mid x^2+y^2 \le 4\}$.

\textbf{Paso 4. Evaluar la integral doble.}
En coordenadas polares, $dA = r\,dr\,d\theta$, con $0 \le r \le 2$,
$0 \le \theta \le 2\pi$:
\[
\iint_R 2\,dA = 2 \int_0^{2\pi} \int_0^2 r\,dr\,d\theta
= 2 \cdot 2\pi \cdot \frac{4}{2} = 8\pi.
\]
Por el teorema de Green, la integral de línea sobre la frontera
orientada positivamente debe ser $8\pi$.

\textbf{Verificación directa.}
Parametrizamos la circunferencia como $\mathbf{r}(t)=\langle 2\cos t,
2\sin t \rangle$, $0 \le t \le 2\pi$. Entonces
\[
\mathbf{F}(\mathbf{r}(t)) = \langle -2\sin t, 2\cos t \rangle,
\quad \mathbf{r}'(t) = \langle -2\sin t, 2\cos t \rangle.
\]
Por tanto,
\[
\oint_C \mathbf{F}\cdot d\mathbf{r}
= \int_0^{2\pi} \langle -2\sin t, 2\cos t \rangle \cdot
\langle -2\sin t, 2\cos t \rangle \, dt
= \int_0^{2\pi} 4(\sin^2 t + \cos^2 t)\,dt
= \int_0^{2\pi} 4\,dt = 8\pi.
\]
La igualdad se cumple: $\boxed{8\pi = 8\pi}$. $\square$
\end{myproof}

%%==========================================================
\section{Divergencia y rotacional}
%%==========================================================

El teorema de Green introduce la cantidad $\partial N/\partial x -
\partial M/\partial y$, que es precisamente la componente $z$ del
rotacional de $\mathbf{F}$ en dos dimensiones. Para campos en
$\mathbb{R}^3$, el rotacional es un vector que mide la tendencia del
campo a \emph{girar} alrededor de un punto, mientras que la divergencia
mide la tendencia a \emph{fluir} hacia el exterior o el interior.

\begin{definition}[Divergencia]
La \textbf{divergencia} de un campo vectorial $\mathbf{F}(x,y,z) =
\langle F_1, F_2, F_3 \rangle$ es el campo escalar
\[
\nabla \cdot \mathbf{F}
= \frac{\partial F_1}{\partial x}
+ \frac{\partial F_2}{\partial y}
+ \frac{\partial F_3}{\partial z}.
\]
En $\mathbb{R}^2$, para $\mathbf{F} = \langle M, N \rangle$,
\[
\nabla \cdot \mathbf{F} = \frac{\partial M}{\partial x}
+ \frac{\partial N}{\partial y}.
\]
\end{definition}

\begin{definition}[Rotacional]
El \textbf{rotacional} de un campo vectorial $\mathbf{F}(x,y,z) =
\langle F_1, F_2, F_3 \rangle$ es el campo vectorial
\[
\nabla \times \mathbf{F}
= \begin{vmatrix}
\mathbf{i} & \mathbf{j} & \mathbf{k} \\
\partial_x & \partial_y & \partial_z \\
F_1 & F_2 & F_3
\end{vmatrix}
= \left( \frac{\partial F_3}{\partial y} - \frac{\partial F_2}{\partial z} \right)
\mathbf{i}
+ \left( \frac{\partial F_1}{\partial z} - \frac{\partial F_3}{\partial x} \right)
\mathbf{j}
+ \left( \frac{\partial F_2}{\partial x} - \frac{\partial F_1}{\partial y} \right)
\mathbf{k}.
\]
En $\mathbb{R}^2$, el rotacional se reduce al escalar (componente $z$)
\[
(\nabla \times \mathbf{F}) \cdot \mathbf{k}
= \frac{\partial N}{\partial x} - \frac{\partial M}{\partial y},
\]
donde $\mathbf{F} = \langle M, N, 0 \rangle$.
\end{definition}

\begin{rem}[Interpretación física]
\begin{itemize}
    \item La \textbf{divergencia} mide la densidad de flujo neto que
    sale de un punto. Si $\nabla\cdot\mathbf{F}>0$, el punto es una
    \emph{fuente}; si es negativo, un \emph{sumidero}. Un campo
    incompresible (como el flujo de un fluido a densidad constante)
    tiene divergencia cero.
    \item El \textbf{rotacional} mide la circulación local del campo
    alrededor de un punto. Su dirección indica el eje de rotación, y
    su magnitud, la intensidad del giro. Un campo con rotacional cero
    se llama \emph{irrotacional}; si además el dominio es simplemente
    conexo, el campo es conservativo.
\end{itemize}
\end{rem}

En el contexto del teorema de Green, el integrando de la integral doble
es precisamente la componente $z$ del rotacional de $\mathbf{F}$:
\[
\nabla \times \mathbf{F} = \left( \frac{\partial N}{\partial x}
- \frac{\partial M}{\partial y} \right) \mathbf{k}.
\]
Por tanto, el teorema de Green se puede reescribir como
\[
\oint_{\partial R} \mathbf{F} \cdot d\mathbf{r}
= \iint_R (\nabla \times \mathbf{F}) \cdot \mathbf{k} \, dA,
\]
que es un caso particular en el plano del teorema de Stokes en
$\mathbb{R}^3$ (que se estudiará en la lección L4.2).

\begin{rem}[Campos conservativos y rotacional]
Un campo $\mathbf{F}$ definido en un dominio simplemente conexo es
conservativo si y solo si $\nabla \times \mathbf{F} = \mathbf{0}$
(en $\mathbb{R}^3$) o $\partial N/\partial x = \partial M/\partial y$
(en $\mathbb{R}^2$). Esta es una reformulación práctica de la condición
de simetría de las derivadas cruzadas.
\end{rem}

\begin{rem}[El teorema de Green como herramienta de cálculo]
El teorema de Green no solo es un resultado teórico: permite convertir
integrales de línea difíciles en integrales dobles más manejables, y
viceversa. También se usa para calcular áreas de regiones planas
eligiendo $\mathbf{F}$ tal que $\partial N/\partial x - \partial
M/\partial y = 1$, por ejemplo $\mathbf{F} = \langle 0, x \rangle$ o
$\mathbf{F} = \langle -y, 0 \rangle$, de modo que
\[
\text{Área}(R) = \oint_{\partial R} x\,dy = -\oint_{\partial R} y\,dx.
\]
Este es el principio del planímetro, un instrumento mecánico para medir
áreas de regiones irregulares.
\end{rem}

%%==========================================================
\section{Problemas propuestos}
%%==========================================================

\begin{prob}
Calcule la integral de línea $\int_C \mathbf{F}\cdot d\mathbf{r}$ para
$\mathbf{F}(x,y)=\langle x^2, y^2 \rangle$ a lo largo de los siguientes
caminos que unen $(0,0)$ con $(1,1)$:
\begin{enumerate}[(a)]
    \item la línea recta $y=x$;
    \item la parábola $x=t$, $y=t^2$;
    \item la curva quebrada que va de $(0,0)$ a $(1,0)$ y luego a $(1,1)$.
\end{enumerate}
Compare los resultados y discuta si el campo es conservativo.
\end{prob}

\begin{prob}
Determine si cada campo es conservativo. En caso afirmativo, halle su
función potencial.
\begin{multicols}{2}
\begin{enumerate}[(a)]
    \item $\mathbf{F}(x,y)=\langle 2x, 2y \rangle$
    \item $\mathbf{F}(x,y)=\langle y^2, 2xy \rangle$
    \item $\mathbf{F}(x,y)=\langle \cos y, -x\sin y \rangle$
    \item $\mathbf{F}(x,y)=\langle e^x \cos y, -e^x \sin y \rangle$
    \item $\mathbf{F}(x,y,z)=\langle 2x, 2y, 2z \rangle$
    \item $\mathbf{F}(x,y,z)=\langle yz, xz, xy \rangle$
\end{enumerate}
\end{multicols}
\end{prob}

\begin{prob}
Utilice el teorema fundamental de las integrales de línea para calcular
$\int_C \mathbf{F}\cdot d\mathbf{r}$ donde $\mathbf{F}(x,y) =
\langle 2xy^2, 2x^2y \rangle$ y $C$ es cualquier curva suave por partes
que une $(0,0)$ con $(1,2)$.
\end{prob}

\begin{prob}
Un campo de fuerzas $\mathbf{F}(x,y) = \langle e^x \cos y + y,
-e^x \sin y + x \rangle$ actúa sobre una partícula.
\begin{enumerate}[(a)]
    \item Demuestre que $\mathbf{F}$ es conservativo.
    \item Encuentre la función potencial $f$.
    \item Calcule el trabajo realizado por $\mathbf{F}$ al mover una
    partícula desde $(0,0)$ hasta $(1,\pi/2)$.
\end{enumerate}
\end{prob}

\begin{prob}
Verifique el teorema de Green para el campo $\mathbf{F}(x,y) =
\langle x^2 y, x y^2 \rangle$ y el rectángulo $0\le x\le 2$,
$0\le y\le 1$.
\end{prob}

\begin{prob}
Use el teorema de Green para calcular
\[
\oint_C (x^2 + y^2)\,dx + (2xy)\,dy,
\]
donde $C$ es el cuadrado con vértices $(0,0)$, $(2,0)$, $(2,2)$,
$(0,2)$ orientado positivamente.
\end{prob}

\begin{prob}
Calcule la integral de línea
\[
\oint_C -y^3\,dx + x^3\,dy
\]
donde $C$ es el círculo $x^2+y^2=4$ orientado positivamente, usando
el teorema de Green. Compare con la evaluación directa.
\end{prob}

\begin{prob}
Sea $R$ una región plana con frontera $C$. Demuestre que el área de $R$
está dada por
\[
A(R) = \frac{1}{2} \oint_C x\,dy - y\,dx.
\]
Use esta fórmula para calcular el área de la elipse $\frac{x^2}{a^2}
+\frac{y^2}{b^2}=1$.
\end{prob}

\begin{prob}
Calcule la divergencia y el rotacional de los siguientes campos:
\begin{multicols}{2}
\begin{enumerate}[(a)]
    \item $\mathbf{F}(x,y,z)=\langle x^2, y^2, z^2 \rangle$
    \item $\mathbf{F}(x,y,z)=\langle yz, xz, xy \rangle$
    \item $\mathbf{F}(x,y,z)=\langle e^x, \sin y, \cos z \rangle$
    \item $\mathbf{F}(x,y,z)=\langle -y, x, 0 \rangle$
\end{enumerate}
\end{multicols}
\end{prob}

\begin{prob}
Demuestre que para cualquier campo $\mathbf{F}$ con segundas derivadas
parciales continuas, se cumple $\nabla \cdot (\nabla \times \mathbf{F})
= 0$. Esto muestra que el rotacional de un campo siempre es
\emph{solenoidal} (sin divergencia). Use un ejemplo concreto para
verificar esta identidad.
\end{prob}

\begin{prob}
Interprete físicamente el resultado del problema anterior:
si $\mathbf{F}$ es el campo de velocidades de un fluido, ¿qué significa
que $\nabla \cdot (\nabla \times \mathbf{F})=0$?
\end{prob}

\begin{prob}
Determine si el campo $\mathbf{F}(x,y) = \langle -y/(x^2+y^2),
x/(x^2+y^2) \rangle$ es conservativo en su dominio. Calcule su integral
de línea alrededor de la circunferencia unidad y explique por qué el
resultado no contradice el teorema fundamental de las integrales de
línea.
\end{prob}